\begin{table}[!htbp]
\centering
\caption{Permanent public-spending increase: stationary effects}
\label{tab:steady_results}
\begin{threeparttable}
\begin{tabular}{lrrrrr}
\toprule
Financing of additional spending & $\Delta K$ & $\Delta Y$ & $\Delta C$ & $\Delta$ informality & $\Delta\tau_f$ \\
 & (\%) & (\%) & (\%) & (p.p.) & (p.p.) \\
\midrule
Lump-sum financing & 0.00 & 0.00 & -2.80 & 0.00 & 0.00 \\
Mixed financing & -8.17 & -3.83 & -6.58 & 3.64 & 3.69 \\
Payroll financing & -22.08 & -10.35 & -13.02 & 9.85 & 9.84 \\
\bottomrule
\end{tabular}
\begin{tablenotes}[flushleft]\footnotesize
\item Notes: The increase in government purchases equals two percent of benchmark GDP. All scenarios share the same initial steady state. Lump-sum financing means that none of the additional spending is charged to the formal wage bill; mixed and payroll financing charge 50 and 100 percent, respectively.
\end{tablenotes}
\end{threeparttable}
\end{table}
